\section{Benchmark Evaluation}\label{sec:experimental_evaluation}

\subsection{Evaluation Framework}\label{sec:experimental_setup}


\paragraph{Evaluation Paradigms.}
Formal Conjectures provides a dynamic, rigorous benchmark for advanced automated reasoning premised on novel mathematical insight as the ultimate intellectual hurdle for complex problems. We define two distinct setups:
\begin{enumerate}
    \item \textbf{Primary Goal: New Formal Proof Discovery.}
    Measuring AI's ability to discover formal proofs for open problems, these unsolved conjectures form a zero-contamination testbed for genuine mathematical discovery. Since no solutions exist in any training corpus, success provides a definitive signal of reasoning capabilities, removing the need for secretive evaluation sets common in other benchmarks.
    \item \textbf{Secondary Goal: Proof Auto-formalization.}
    This track provides a climbable benchmark treating the non-trivial translation of established mathematics into formal code as a distinct, rigorous challenge. It measures a model's ability to work with Lean~4 and Mathlib to formalize known arguments.
\end{enumerate}


\paragraph{Evaluating Open Problems.}
Evaluating open problems is definitive: as no formal solutions exist at release time, the first kernel-accepted proof provides a zero-contamination success signal.
However, since a proof may settle a misformalized statement rather than the intended problem, any result for an open conjecture triggers manual inspection for fidelity.


\paragraph{Fixed Benchmark Subsets.}
To enable stable model comparisons, we provide two frozen subsets of 100 problems each:
\texttt{FC100OpenSet1} (100 \texttt{research open} statements) and \texttt{FC100SolvedSet1} (100 \texttt{research solved} statements).
These subsets are defined in the repository files \texttt{FC100OpenSet1.lean} and \texttt{FC100SolvedSet1.lean}, which import exactly the corresponding theorem statements.
Because these files are compiled as part of the repository across all supported Lean versions, the problem sets remain fixed and comparable even as the repository evolves.
Further details on the construction of these subsets are in Appendix~\ref{sec:app_bench_lists}.


\paragraph{Correctness.}
Following established formal math benchmark methodologies \citep{minif2f21, putnambench24}, our evaluation leverages the Lean~4 kernel for an objective, rigorous, and automatable success criterion.
A proposed solution is a Lean proof term replacing the \lstinline{sorry} placeholder.
Solutions are correct if and only if accepted by the kernel without relying on forbidden axioms, such as \lstinline{sorry}.
This binary criterion eliminates human grading ambiguity and provides an indisputable ground truth.
While theoretically susceptible to foundational bugs in the Lean kernel or unintended axioms (e.g., a Lean Zulip thread\footnote{\url{https://leanprover.zulipchat.com/\#narrow/channel/270676-lean4/topic/Soundness.20bug.3A.20hasLooseBVars.20is.20not.20conservative/with/520153084}}, which mentions a now-fixed exploit), in practice, kernel-level verification provides the highest mathematical rigor.
Consequently, all evaluations are automatically verifiable, allowing results and proof terms to be shared publicly for transparent, reproducible research.


\paragraph{Versioning and Reproducibility.}\label{sec:versioning}
Evaluating on a living repository is challenging: dependency updates can alter statement semantics, and new Mathlib theorems may simplify problems.
To enable reproducible comparisons, we tag frozen benchmark snapshots using a two-part naming scheme: \texttt{bench-v$N$-lean4.$X$.$Y$}.
The \emph{bench version} \texttt{v$N$} identifies a fixed problem set; the \emph{Lean version} suffix pins the Mathlib tag (and hence Lean toolchain) against which statements compile and are evaluated.
Releases are issued every few months; when Lean versions bump, companion tags allow evaluating the same problem set against intermediate toolchains, e.g.\ \texttt{bench-v3-lean4.27.0}.
To preserve baselines, snapshots are immutable; misformalization fixes are never patched into existing versions but incorporated into the next (e.g., \texttt{bench-v$($N$+1)$}).


\paragraph{A Living Benchmark and Unified API.}
Outside the frozen subsets, the repository operates as a dynamic benchmark that exhibits useful evaluation dynamics over time.
As new conjectures are added, solved, or corrected by the community, it naturally self-adjusts its difficulty, expands, and improves.
The hardest problems remain open, pushing the frontier of automated reasoning.
Crucially, the repository acts as a unified API: researchers can submit formalized conjectures to a single centralized hub for simultaneous exposure to all evaluating AI systems.
This obviates the need to manually engage with individual provers, ensuring broad and concurrent attempts at mathematical discovery.

\begin{table}[t!]
  \caption{Proof coverage across all statement categories. \emph{Proved in repo}: the statement has a \lstinline{sorry}-free proof in the repository; \emph{Linked proof}: a proof is recorded via the \lstinline{@[formal\_proof]} attribute but is not present in the repository itself; \emph{With proof}: either of the above.}
  \label{tab:formal_proof_baseline}
  \centering
  \begin{tabular}{lrrrr}
    \toprule
    Category & Total & \multicolumn{3}{c}{Proof Status} \\
    \cmidrule(lr){3-5}
    & & Proved in repo & Linked proof & With proof \\
    \midrule
    \texttt{research\_solved} & 836 & 44 (5.3\%) & 101 (12.1\%) & 145 (17.3\%) \\
    \texttt{textbook} & 128 & 45 (35.2\%) & 3 (2.3\%) & 48 (37.5\%) \\
    \texttt{test} & 467 & 390 (83.5\%) & 3 (0.6\%) & 393 (84.2\%) \\
    \texttt{API} & 155 & 140 (90.3\%) & 0 & 140 (90.3\%) \\
    \midrule
    \textbf{Total} & 1586 & 619 (39.0\%) & 107 (6.7\%) & 726 (45.8\%) \\
    \bottomrule
  \end{tabular}
\end{table}



\paragraph{Open Problems Baseline.}
For the primary proof-discovery task on open problems, the baseline for the first tagged release \texttt{bench-v1-lean4.27} is $0\%$ for all systems: by definition, no formal proof exists for any \texttt{research open} statement at the time of tagging (problems with known solutions are reclassified as \texttt{research solved}).
Over time, these open statements may receive formal or informal proofs.
Proofs that expose a misformalization (Section~\ref{sec:misform}) rather than settling the intended problem trigger revisions incorporated into the next version.
In subsequent releases, interim-solved problems are reclassified to \texttt{research solved}, with formal solutions receiving the \lstinline{@[formal_proof]} attribute.
The new \texttt{research open} set combines remaining unsolved problems with new additions; thus, successive versions start with a baseline of $\geq 0\%$ solved statements and become increasingly more difficult.

\paragraph{Proof Auto-Formalization Baseline.}
For auto-formalization, the baseline is system-agnostic: we track all formal proofs contributed to the repository, regardless of origin.
These are tracked using the \lstinline{@[formal_proof]} attribute (Section~\ref{sec:formalization-methodology}), which records the proof system (\texttt{formal\_conjectures} for proofs within the repository, \texttt{lean4} for proofs in external Lean~4 projects such as Mathlib, or \texttt{other\_system} for proofs in Isabelle, Rocq, etc.) and links to the source.
Table~\ref{tab:formal_proof_baseline} summarizes the current state. Notably, 17.3\% of \texttt{research\_solved} statements have a known formal proof (5.3\% internal, 12.1\% linked externally).


\subsection{Illustrative Evaluation}\label{sec:experimental_results}


We provide illustrative evaluations, with setup and cost details in Appendix~\ref{sec:illustrative_evaluation}.


\begin{wraptable}{r}{0.45\textwidth}
  \vspace{-15pt}
  \caption{Results on the frozen \texttt{FC100SolvedSet1}.}
  \vspace{5pt}
  \label{tab:results_bench_table}
  \centering
  \small
  \begin{tabular}{lr}
    \toprule
    Method & Proved (\%) \\
    \midrule
    AlphaProof (1k sims)  & 45.0\% \\
    AlphaProof (16k sims) & 50.0\% \\
    DM prover agent (dev)       & 66.0\% \\
    \bottomrule
  \end{tabular}
  \vspace{-10pt}
\end{wraptable}


\paragraph{Frozen Evaluation Subsets.}
To ensure stable comparisons despite repository fluidity, we provide two frozen evaluation subsets, each containing 100 problems randomly sampled from the \texttt{bench-v1-lean4.27.0} tag: \texttt{FC100SolvedSet1} (proof auto-formalization) and \texttt{FC100OpenSet1} (proof discovery; 0\% baseline).
By definition, \texttt{FC100OpenSet1} currently remains at a definitive 0\% baseline across these evaluated methods, which highlights its status as a rigorous frontier for new discovery.
Results on \texttt{FC100SolvedSet1} (Table \ref{tab:results_bench_table}) demonstrate a clear, climbable signal across both compute and model capabilities.
To provide a diverse evaluation, we evaluate both a slightly improved version of the AlphaProof system used in the 2024 IMO \citep{AlphaProof25} and a development version of a DeepMind prover agent \citep{DMProverAgent26}.
These are different systems rather than incremental versions of the same architecture.
Utilizing a constrained tree-search-only inference setup on v6e TPUs, AlphaProof achieves a solve rate of 45\% on \texttt{FC100SolvedSet1} with a low compute configuration of 1,000 simulations per problem.
By increasing the compute budget to 16,000 simulations, the solve rate improves to 50\%, demonstrating a clear scaling signal.
We emphasize that these evaluations are conducted using tree-search inference only and do not utilize AlphaProof's test-time reinforcement learning (TTRL) mechanism, which generates a curriculum to learn from during the proof finding process through weight updates.
Furthermore, utilizing a newer, development version of the DeepMind prover agent, the solve rate increases to 66\%, illustrating that the benchmark effectively measures advancements across improved model architectures and methodologies.
Because running automated reasoning methods within Lean involves significant infrastructure overhead, we provide these specific systems as illustrative baselines to validate the benchmark's signal.
We encourage the community to report results against these subsets as well as the evolving main repository.
